\section[The node tree: PartNode and PartsModel]{The node tree: \code{PartNode} and
\code{PartsModel}}\label{sec:tree}

Where the leaf of \cref{sec:leaf} is the physical vocabulary, the tree is the grammar. Two classes
in \srcf{core/model/PartsModel.h} carry all of it. \code{PartNode} is one interior node of the
rocket: it owns exactly one \code{Part}, the \code{StationLink} that seats it on its parent
(\cref{sec:placement}), and the child nodes beneath it --- so any node is a complete sub-assembly,
and the composite queries of \cref{sec:caching} run identically on a full rocket or a booster
stack. \code{PartsModel} is the single owner of that tree and its only mutation authority: reads
hand out \cppinline|const Part&|, and every post-attach change --- structural or leaf-level --- is
a named verb on the model that performs the write, maintains the id index and the motor borrow,
dirties the correct cache chain, and brackets itself in change events. \Cref{fig:tree-class} shows
the full ownership picture, with \code{RocketModel} above as the value owner and install facade.

This section covers the tree's structure and mutation surface: the node's owned interior and read
surface, how detached trees are built and cloned, lookup and traversal, the structural verbs and
their typed errors, the routed leaf edits with their invalidation hint, and the event stream that
lets a Qt front-end mirror the tree without the model knowing Qt exists. The caches that hang off
each node are specified in \cref{sec:caching}; the motor slot's semantics in \cref{sec:motor}.

\classfig{class-partsmodel}{The owned part tree (\srcfcap{core/model/PartsModel.h}).
\code{RocketModel} holds \code{PartsModel} by value; \code{PartsModel} owns the root node through
\cppinline|unique_ptr| and indexes every node by part id; each \code{PartNode} owns its leaf and
its children the same way. \code{Event} and the error enums are nested types of
\code{PartsModel}.}{fig:tree-class}{0.85}

\subsection{The node: an owned interior, a const read surface}\label{sec:tree-node}

A \code{PartNode} holds five members (\src{core/model/PartsModel.h}{108}--112). Three express
ownership and shape: \cppinline|part_|, a \cppinline|std::unique_ptr<part::Part>| --- the node
\emph{is} the owner of its leaf; \cppinline|children_|, a
\cppinline|std::vector<std::unique_ptr<PartNode>>| in attachment order; and \cppinline|link_|, the
\code{StationLink} stored on the \emph{child} as its incoming edge --- the placement intent
relative to whatever parent the node currently has, a default link on a root. Storing the edge on
the child rather than in a parent-side table means a detached sub-tree carries its internal
seating with it: every edge except the one that was cut travels inside the value.

The remaining two are raw back-pointers, and both are deliberate non-owners.
\cppinline|parent_| points up the tree; it is safe as a bare pointer because ownership already
guarantees the parent outlives the child --- destroying a parent destroys its
\cppinline|children_| first. \cppinline|model_| points at the owning \code{PartsModel} while the
node is attached to one and is null while detached; it exists purely as a guard bit for
\code{addChild} (\cref{sec:tree-detached}). Because nodes point back at their model,
\code{PartsModel} is non-copyable and non-movable (\src{core/model/PartsModel.h}{156}) --- a
moved-from model would leave every node's back-pointer dangling, so the type simply forbids it.

The read surface is small and entirely \code{const}: \code{id()} forwards to the part's own id, so
there is exactly one identity namespace --- a node is named by the part it owns, and no separate
node-id can ever disagree with it; \code{part()} returns \cppinline|const Part&|, so no reader can
mutate a leaf behind the model's back; \code{link()}, \code{parent()}, and \code{children()} (a
\cppinline|std::span| view, exposing the sequence but not the vector) complete the topology; and
\code{rowInParent()} reports the node's index among its siblings, 0 for a root
(\src{core/model/PartsModel.cpp}{55}). That last accessor earns its place by shape: (parent, row)
is precisely the coordinate system of \code{QAbstractItemModel}, so the Qt adapter of
\cref{sec:gui} reads its geometry straight off the node instead of maintaining a parallel map.

\subsection[Building detached trees: make, addChild, clone]{Building detached trees: \code{make}, \code{addChild}, \code{clone}}
\label{sec:tree-detached}

Nodes begin life detached. The only way to create one is the static factory
\code{PartNode::make(p, link)} (\src{core/model/PartsModel.h}{47}), which wraps a leaf and its
intended incoming edge into a fresh node --- null part in, null node out. A detached tree is then
assembled bottom-up with \code{addChild}, which wires the child's \cppinline|parent_| pointer,
stores the edge, and moves the child in (\cref{lst:tree-addchild}). This is the loader's idiom:
\code{DesignSerializer} builds the whole tree from the file --- \code{make} per part,
\code{addChild} per edge (\src{core/model/DesignSerializer.cpp}{163}) --- and only then hands the
finished root to the model in one step.

\begin{listing}[htbp]
\begin{minted}{cpp}
void PartNode::addChild(std::unique_ptr<PartNode> child, part::StationLink link)
{
    if(!child)
    {
        utils::Logger::getInstance()->error("PartNode::addChild: ignoring null child");
        return;
    }
    if(model_ != nullptr || child->model_ != nullptr)
    {
        utils::Logger::getInstance()->error(
            "PartNode::addChild: owned trees mutate only through PartsModel verbs");
        return;
    }
    child->parent_ = this;
    child->link_   = link;
    children_.push_back(std::move(child));
    markMassDirty();
    markPlacementDirty();
}
\end{minted}
\caption{Detached-tree assembly and its ownership guard
(\srccap{core/model/PartsModel.cpp}{72}--90).}\label{lst:tree-addchild}
\end{listing}

The guard in the middle of \cref{lst:tree-addchild} is the boundary of the whole design:
\code{addChild} is a logged no-op the moment either node is owned by a model.

\begin{designrule}[Owned trees mutate only through verbs]
Once a node belongs to a \code{PartsModel}, a structural edit implies bookkeeping the node cannot
do alone: the id index must gain or lose entries, the motor borrow may need re-resolving, cache
chains must be dirtied \emph{from the right node upward}, and the Qt-facing event pair must fire
with correct pre-state. \code{addChild} on an owned node would silently skip all four, so it
refuses --- the \cppinline|model_| back-pointer exists to make ``am I owned?'' answerable at the
one place it matters. Detached trees have none of those obligations, which is exactly why free
assembly is legal there.
\end{designrule}

\code{clone()} (\src{core/model/PartsModel.cpp}{92}) deep-copies a sub-tree by recursion on this
same machinery: every leaf duplicated through the virtual \code{Part::clone} --- so each copy
carries a fresh id, per the leaf's identity contract (\cref{sec:leaf}) --- and every
\code{StationLink} copied verbatim, so the copy is seated exactly as the original was. The result
is a detached root, ready for \code{attachSubtree}: clone-then-attach is copy-paste.

\subsection{The model: identity, lookup, traversal}\label{sec:tree-lookup}

\code{PartsModel} holds a \emph{nullable} root (\cppinline|std::unique_ptr<PartNode> root_|) and
reports it through \code{hasDesign()} (\src{core/model/PartsModel.h}{161}). The null state is a
real state, not a placeholder: the GUI opens with no design loaded, the CLI runs before any
\code{loaddesign}, and both must render, save-refuse, and report sensibly. Representing ``no
design'' as an empty-but-present dummy root would force every consumer to distinguish a real
single-part rocket from the sentinel; a null root makes the distinction structural, and consumers
gate on \code{hasDesign()} before asking composite questions (the CLI's design commands and the
geometry-derived reference area both check it explicitly, \cref{sec:cli,sec:flight}).

Lookup is by id, in constant time: \cppinline|index_| maps
\cppinline|Part::Id| $\to$ \cppinline|PartNode*| for every owned node, and \code{find(id)} returns
the node or \cppinline|nullptr| (\src{core/model/PartsModel.h}{167}). The index is maintained
exclusively by the adoption pair \code{registerSubtree}/\code{unregisterSubtree}
(\cref{sec:tree-verbs}), so it can never disagree with the ownership tree.

\begin{keyinsight}[Stale ids fail safe]
Front-ends hold ids across time: a GUI selection survives a delete, a CLI script names parts that
a previous command may have removed, an undo stack refers to nodes that no longer exist. Every
verb therefore resolves its id through \code{find} and returns a typed error or \code{false} on a
miss --- and because ids are drawn from a process-wide counter and never reused within a run
(\cref{sec:leaf}), a stale id can never alias a \emph{different} part later. The worst outcome of
a dangling identity is a refused verb with a diagnosable error, never a write to the wrong node
and never undefined behavior. This is the payoff of making the id, not a pointer or a name, the
currency of the mutation API.
\end{keyinsight}

\code{findByName} (\src{core/model/PartsModel.cpp}{257}) is the deliberate second-class citizen:
a pre-order first-match walk, provided because humans type names at the CLI, but names carry no
uniqueness requirement --- two \code{"BodyTube"}s are legal --- so the method is a convenience
for interactive lookup and nothing in the model keys on it. \code{size()} is the index's size,
$O(1)$ and zero with no design. Traversal is \code{forEachNode}
(\src{core/model/PartsModel.h}{185}), a template over \cppinline|f(const PartNode&, int depth)|
that visits parent-before-child in attachment order --- the same deterministic order the placement
resolver emits (\cref{sec:placement}), so a listing printed from the walk and a diagnostic
reported from the resolver name parts in the same sequence.

\subsection{Structural verbs and typed errors}\label{sec:tree-verbs}

Four verbs change the shape of an owned tree, and each returns its failure modes in the type:
\code{attach} and \code{attachSubtree} return
\cppinline|std::expected<Part::Id, AttachError>|, \code{detach} returns
\cppinline|std::expected<std::unique_ptr<PartNode>, DetachError>|, and \code{installRoot} is
total. The enums are small and exact --- \code{AttachError} is
\{\code{NullPart}, \code{NoSuchParent}, \code{DuplicateMotor}\}, \code{DetachError} is
\{\code{NoSuchId}, \code{IsRoot}\} --- so a front-end can map each case to a specific message
instead of parsing a boolean (\cref{sec:cli}). \Cref{fig:tree-attach} traces the attach path;
\cref{lst:tree-attach} is its body.

\tikzfig{activity-attach}{The attach path
(\srcfcap{core/model/PartsModel.cpp}). \code{attach} wraps the bare part and delegates to
\code{attachSubtree}; every failure exits before any state changes, and the event pair brackets
exactly the mutating region.}{fig:tree-attach}

\begin{listing}[htbp]
\begin{minted}{cpp}
std::expected<part::Part::Id, PartsModel::AttachError>
PartsModel::attachSubtree(part::Part::Id parent, std::unique_ptr<PartNode> subtree,
                          std::optional<part::StationLink> linkOverride)
{
    if(!subtree || subtree->model_ != nullptr)
    {
        return std::unexpected(AttachError::NullPart);
    }
    PartNode* parentNode = find(parent);
    if(parentNode == nullptr)
    {
        return std::unexpected(AttachError::NoSuchParent);
    }
    // Single-motor invariant, enforced where it can be broken: reject a second Motor
    // anywhere in the incoming sub-tree.
    if(motor_ != nullptr)
    {
        // ... pre-order scan of the incoming sub-tree ...
        if(hasMotor)
        {
            return std::unexpected(AttachError::DuplicateMotor);
        }
    }

    if(linkOverride)
    {
        subtree->link_ = *linkOverride;
    }
    const part::Part::Id childId = subtree->id();
    const int row = static_cast<int>(parentNode->children_.size());
    const Event e{Event::Attached, childId, parentNode->id(), row};

    emitEvent(e, true);
    subtree->parent_ = parentNode;
    registerSubtree(subtree.get());
    parentNode->children_.push_back(std::move(subtree));
    parentNode->markMassDirty();
    parentNode->markPlacementDirty();
    emitEvent(e, false);
    return childId;
}
\end{minted}
\caption{\code{attachSubtree}, lightly elided
(\srccap{core/model/PartsModel.cpp}{304}--350): guards first, then the event-bracketed adoption.}
\label{lst:tree-attach}
\end{listing}

Three details of \cref{lst:tree-attach} carry most of the design. First, \code{attach} is
literally \code{attachSubtree} over a single-node tree
(\src{core/model/PartsModel.cpp}{298}) --- one adoption path, not two. Second, the single-motor
invariant is checked \emph{here}, at the only door a second motor could enter through: when a
motor is already installed, the incoming sub-tree is scanned and refused with
\code{DuplicateMotor} if it contains one (\cref{sec:motor} covers why at most one). Third,
\code{linkOverride} exists for re-seating: a sub-tree built by the serializer arrives with its
stored links and should keep them, but a pasted or moved sub-tree is being seated on a \emph{new}
parent, and its old incoming edge is intent about a parent it no longer has. The override replaces
exactly that one edge --- the cut edge --- and leaves the sub-tree's interior seating untouched.

\code{detach(id)} is the inverse (\src{core/model/PartsModel.cpp}{352}), and its return type is
the point: the removed sub-tree comes back as a live
\cppinline|std::unique_ptr<PartNode>| --- unregistered, parent pointer cleared, internal links
intact. A first-class detached value is the building block for every higher-level edit:
detach-then-drop is delete (the CLI's \code{removepart}, \src{cli/Repl.cpp}{1056}),
detach-then-attachSubtree is move, hold-and-reattach is one undo frame, and clone-then-attach is
paste. None of those features needs model support beyond the verb pair. The root is the one node
\code{detach} refuses (\code{IsRoot}): it is the datum every resolved pose is expressed against,
and removing it is not an edit to the tree but a replacement of the design ---
which is what \code{installRoot} is for. \code{installRoot} atomically swaps the whole tree
(null clears), unregistering the old root and registering the new
(\src{core/model/PartsModel.cpp}{380}).

All four verbs converge on the private adoption pair.
\code{registerSubtree}/\code{unregisterSubtree} (\src{core/model/PartsModel.cpp}{270}) walk the
sub-tree once and maintain the three pieces of derived ownership state together: the id index
gains or loses every node, each node's \cppinline|model_| back-pointer is set or cleared, and the
typed motor borrow (\cppinline|motor_|) is re-resolved if the walk crosses a \code{Motor}. That
pair is the \emph{only} code that touches \cppinline|motor_| --- attach, detach, and a clearing
\code{installRoot} all route through it, so a cleared design can never leave the borrow dangling
(\cref{sec:motor}).

\subsection{Routed leaf edits and the mutation hint}\label{sec:tree-edits}

Post-attach leaf edits reach the tree through four routed verbs, the public face of the
private-plus-friend seam on \code{Part} (\cref{sec:leaf}). \code{setPartMass} and
\code{setPartInertia} (\src{core/model/PartsModel.cpp}{399}) perform the plain write and dirty the
mass chain from the edited node upward --- and only the mass chain. \code{setLink}
(\src{core/model/PartsModel.cpp}{431}) re-seats an existing edge (station, gap, or seat-kind
edits; refused on a root, which has no incoming edge) and dirties \emph{both} chains from the
parent: a seat edit moves resolved poses, and an equal-mass geometry change is exactly what the
mass-delta gate of \cref{sec:caching} cannot see, so relying on the mass gate alone would serve a
stale CM.

\code{mutatePart} (\cref{lst:tree-mutate}) is the escape hatch for edits with no typed verb: an
arbitrary functor applied to the leaf under a declared \code{MutateHint}.

\begin{listing}[htbp]
\begin{minted}{cpp}
bool PartsModel::mutatePart(part::Part::Id id, MutateHint hint,
                            const std::function<void(part::Part&)>& fn)
{
    PartNode* n = find(id);
    if(n == nullptr || !fn)
    {
        return false;
    }
    const Event e{Event::Mutated, id, n->parent_ ? n->parent_->id() : part::Part::Id{0},
                       n->rowInParent()};
    emitEvent(e, true);
    fn(*n->part_);
    n->markMassDirty();
    if(hint == MutateHint::Geometry)
    {
        n->markPlacementDirty();
    }
    emitEvent(e, false);
    return true;
}
\end{minted}
\caption{The hinted escape hatch (\srccap{core/model/PartsModel.cpp}{449}--468): the write and its
invalidation are fused, and the hint chooses how much invalidation.}\label{lst:tree-mutate}
\end{listing}

\begin{designrule}[Declare what an edit can move]
The two invalidation chains have very different prices. Dirtying the mass chain schedules an
$O(N)$ re-weighting at the next composite read; dirtying the placement chain schedules a full
geometry re-resolve plus the $O(N \log N)$ envelope sweep (the sort at
\src{core/model/parts/Placement.cpp}{149}). A mass-only edit --- a motor swap, a payload trim ---
moves no pose and cannot create an overlap, so \code{MassOnly} skips the sweep entirely; a length
or radius edit can do both, so \code{Geometry} dirties both chains. The hint is a declaration the
caller must make, not an inference the model attempts: the model cannot see inside the functor,
and guessing conservatively on every edit would put the sweep on paths that run per-swap in batch
motor studies. Note the asymmetry is one-way --- \code{Geometry} still dirties mass, because
reshaping a part moves its CM even at constant mass.
\end{designrule}

\subsection{The event stream}\label{sec:tree-events}

Every mutating verb in the previous two subsections brackets itself in a pair of notifications:
one \code{Event} value fired with \cppinline|before == true| immediately \emph{before} the first
write, and the identical value fired with \cppinline|before == false| after the last
(\cref{lst:tree-event}). The event carries the affected node's id, its (ex-)parent's id, and its
row within that parent --- captured \emph{pre-mutation}, so a \code{Detached} event names the row
the node still occupies and an \code{Attached} event names the row the node is about to land in.

\begin{listing}[htbp]
\begin{minted}{cpp}
/// One change notification, fired as aboutTo/did pairs (before == true, then false)
/// around every mutation -- the shape Qt's begin*/end* row-op contract needs.
struct Event
{
    enum Kind : std::uint8_t { Reset, Attached, Detached, LinkChanged, Mutated };
    Kind kind{Reset};
    part::Part::Id id{0};        ///< the affected node (0 on a clearing Reset)
    part::Part::Id parentId{0};  ///< its (ex-)parent; 0 for a root
    int row{0};                  ///< its row within that parent
};
using ChangeCallback = std::function<void(const Event&, bool before)>;
\end{minted}
\caption{The change notification (\srccap{core/model/PartsModel.h}{143}--151).}
\label{lst:tree-event}
\end{listing}

The shape is not incidental: it is the exact contract of Qt's model/view row operations, which
demand \code{beginInsertRows(parent, row, row)} before the underlying insert and
\code{endInsertRows()} after, with removal mirrored. The GUI adapter is therefore a
\cppinline|switch| that forwards each pair verbatim (\src{gui/RocketTreeView.cpp}{32});
\cref{sec:gui} covers it. The model keeps a single callback slot
(\code{setChangedCallback}; the latest registration wins, \cppinline|{}| clears), and
\code{RocketModel} is that single subscriber: its constructor installs a bridge
(\src{core/model/RocketModel.cpp}{16}) that forwards the typed stream to whichever granular
consumer registered via \cppinline|setPartsEventCallback| and folds each completed mutation
(\cppinline|before == false|) into the coarse ``something changed'' callback that simpler
consumers prefer. One producer, one bridge, two subscription shapes.

\begin{rejected}[Rejected --- Qt signals in the model]
The natural Qt idiom --- make \code{PartsModel} a \code{QObject} and emit signals --- would drag
QtCore into \cppinline|qtrocket_core|, and with it into the CLI and every test binary, breaking the
layering rule that only \srcf{gui/} may depend on Qt. A \cppinline|std::function| callback
carrying aboutTo/did pairs delivers everything the Qt adapter needs while costing the headless
consumers nothing: the CLI simply never registers one. Multicast is deliberately absent ---
\code{RocketModel} is by construction the one subscriber, and it already fans out --- so the model
avoids owning listener-list bookkeeping it would never exercise.
\end{rejected}

\subsection{The install facade}\label{sec:tree-install}

Whole-design replacement has one public door: \code{RocketModel::installDesign}
(\src{core/model/RocketModel.h}{87}), with \code{clearDesign()} as sugar for installing null.
It delegates the tree swap to \code{PartsModel::installRoot} and then performs the one install
side effect that belongs to \code{RocketModel} rather than the tree: it clears the manual
reference-area override (\src{core/model/RocketModel.cpp}{118}). The override is drag-model state,
not tree state --- a provenance flag marking the stored frontal area as user-entered
(\cref{sec:flight}) --- and it describes the airframe it was entered for. A freshly installed
design must not present it as authored intent for the new airframe; resetting the flag strips
that provenance, so a subsequent save records the area as unauthored and a loaded file's own
\code{<sim>} block decides whether an override is re-established (\cref{sec:serialization}).
Both loaders honor the facade: the serializer installs a loaded
design through it (\src{core/model/DesignSerializer.cpp}{262}), and the CLI's design commands
build a detached root and hand it over the same way (\src{cli/Repl.cpp}{880}).

The tree, then, is a closed system: one owner, one id currency, verbs that fuse writes with their
bookkeeping, and an event stream shaped for the one consumer that needs granularity. What those
verbs actually invalidate --- and what the flight loop pays when it reads the composites back ---
is the subject of \cref{sec:caching}.
